% cite as https://ling.auf.net/lingbuzz/009369

% !TEX TS-program = xelatex
\documentclass[12pt,letterpaper]{article}

% !TEX TS-program = xelatex
% File: preamble.tex
% Purpose: LaTeX preamble for this paper
% Version: 2.0.0 (December 2025 typography redesign)
%
% Usage: % !TEX TS-program = xelatex
% File: preamble.tex
% Purpose: LaTeX preamble for this paper
% Version: 2.0.0 (December 2025 typography redesign)
%
% Usage: % !TEX TS-program = xelatex
% File: preamble.tex
% Purpose: LaTeX preamble for this paper
% Version: 2.0.0 (December 2025 typography redesign)
%
% Usage: \input{preamble.tex} in main document
%
% Compiler: XeLaTeX (default) or LuaLaTeX
%           NOT pdfLaTeX - requires fontspec for fonts

% ===========================
% FONTS
% ===========================
\usepackage{fontspec}

\setmainfont{EB Garamond}[
  Numbers=OldStyle,
  Ligatures=TeX,
  BoldFont={EB Garamond},
]

% IPA fallback font
\newfontfamily\ipafont{Charis SIL}
\newcommand{\ipa}[1]{{\ipafont #1}}

% CJK fallback font (Japanese / Chinese / Korean)
\newfontfamily\cjkfont{Hiragino Mincho ProN}
\newcommand{\cjk}[1]{{\cjkfont #1}}

% Lining figures when needed (tables, years in isolation)
\providecommand{\liningnums}[1]{{\addfontfeatures{Numbers=Lining}#1}}

% Monospace
\setmonofont{Inconsolata}[Scale=MatchLowercase]

% ===========================
% PAGE LAYOUT
% ===========================
\usepackage[
  letterpaper,
  inner=1.25in,
  outer=1in,
  top=1in,
  bottom=1.25in,
  marginparwidth=0.6in,
]{geometry}

\usepackage[british]{babel}
\usepackage[final]{microtype}

% ===========================
% HEADINGS
% ===========================
\usepackage{titlesec}

% Section: small caps, number in margin
\titleformat{\section}
  {\normalfont\scshape}
  {\llap{\thesection\quad}}
  {0pt}
  {}

% Subsection: small caps sentence case
\titleformat{\subsection}
  {\normalfont\normalsize\scshape}
  {\thesubsection\quad}
  {0pt}
  {}

% Subsubsection: upright
\titleformat{\subsubsection}
  {\normalfont}
  {\thesubsubsection\quad}
  {0pt}
  {}

% Spacing
\titlespacing*{\section}{0pt}{2ex plus 1ex minus .2ex}{1ex plus .2ex}
\titlespacing*{\subsection}{0pt}{1.5ex plus 1ex minus .2ex}{0.5ex plus .2ex}
\titlespacing*{\subsubsection}{0pt}{1ex plus 0.5ex minus .1ex}{0.3ex plus .1ex}

% ===========================
% RUNNING HEADS
% ===========================
\usepackage{fancyhdr}
\pagestyle{fancy}
\fancyhf{}
\fancyhead[LE]{\small\scshape\leftmark}
\fancyhead[RO]{\small\scshape\@title}
\fancyfoot[LE,RO]{\thepage}
\renewcommand{\headrulewidth}{0pt}

% For oneside documents (most papers):
\fancyhead[L]{\small\scshape\leftmark}
\fancyhead[R]{\small\thepage}
\fancyfoot{}

% ===========================
% COLORS & HYPERLINKS
% ===========================
\usepackage{xcolor}
\definecolor{linkmaroon}{RGB}{128, 0, 32}

\usepackage{hyperref}
\hypersetup{
  colorlinks=true,
  linkcolor=linkmaroon,
  citecolor=linkmaroon,
  urlcolor=linkmaroon,
  pdfauthor={Brett Reynolds},
}

\usepackage{orcidlink}

% ===========================
% QUOTATIONS
% ===========================
\usepackage[style=american]{csquotes} % Double quotes outer (Oxford spelling kept via babel)
\MakeOuterQuote{"}

% ===========================
% SEMANTIC MACROS
% ===========================
% Terms (linguistic concepts being introduced/defined)
\newcommand{\term}[1]{\textsc{#1}}

% Mentions (words/expressions under discussion)
\newcommand{\mention}[1]{\textit{#1}}

% Mentions in headings (angle brackets instead of italics)
% Robust: headings are moving arguments (written to .toc/.aux)
\DeclareRobustCommand{\mentionhead}[1]{⟨\textup{#1}⟩}

% Object language (cited forms, foreign words)
\newcommand{\olang}[1]{\textit{#1}}

% Small-caps abbreviations for glosses
\newcommand{\abbr}[1]{\textsc{#1}}

% Cross-linguistic subscript marker
\usepackage{marvosym}
\newcommand{\crossmark}{\textsubscript{\Cross}}

% ===========================
% MATHS AND SYMBOLS
% ===========================
\usepackage{amsmath,amssymb}

% ===========================
% LINGUISTIC EXAMPLES
% ===========================
\usepackage{langsci-gb4e}
\makeatletter
\@ifundefined{noautomath}{}{\noautomath}
\makeatother

% Judgement markers
\newcommand{\ungram}[1]{*\!#1}
\newcommand{\marg}[1]{?\!#1}
\newcommand{\odd}[1]{\#\!#1}

% ===========================
% LISTS
% ===========================
\usepackage{enumitem}
\setlist{nosep, leftmargin=*}
\setlist[enumerate]{label=\arabic*.}
\setlist[itemize]{label=--}

% ===========================
% TABLES & FIGURES
% ===========================
\usepackage{booktabs}
\usepackage{array}
\usepackage{graphicx}
\graphicspath{{figures/}}

% ===========================
% BIBLIOGRAPHY
% ===========================
\usepackage[backend=biber,style=apa,natbib=true,doi=true,isbn=false,url=true]{biblatex}
\IfFileExists{references-standalone.bib}
  {\addbibresource{references-standalone.bib}}
  {\addbibresource{references.bib}\IfFileExists{references-local.bib}{\addbibresource{references-local.bib}}{}}

% Possessive citation: Author's (Year)
\newcommand{\posscite}[1]{\citeauthor{#1}'s (\citeyear{#1})}

% ===========================
% UTILITIES
% ===========================
\usepackage{xspace}
\newcommand{\eg}{e.g.,\xspace}
\newcommand{\ie}{i.e.,\xspace}
\newcommand{\etc}{etc.\xspace}
 in main document
%
% Compiler: XeLaTeX (default) or LuaLaTeX
%           NOT pdfLaTeX - requires fontspec for fonts

% ===========================
% FONTS
% ===========================
\usepackage{fontspec}

\setmainfont{EB Garamond}[
  Numbers=OldStyle,
  Ligatures=TeX,
  BoldFont={EB Garamond},
]

% IPA fallback font
\newfontfamily\ipafont{Charis SIL}
\newcommand{\ipa}[1]{{\ipafont #1}}

% CJK fallback font (Japanese / Chinese / Korean)
\newfontfamily\cjkfont{Hiragino Mincho ProN}
\newcommand{\cjk}[1]{{\cjkfont #1}}

% Lining figures when needed (tables, years in isolation)
\providecommand{\liningnums}[1]{{\addfontfeatures{Numbers=Lining}#1}}

% Monospace
\setmonofont{Inconsolata}[Scale=MatchLowercase]

% ===========================
% PAGE LAYOUT
% ===========================
\usepackage[
  letterpaper,
  inner=1.25in,
  outer=1in,
  top=1in,
  bottom=1.25in,
  marginparwidth=0.6in,
]{geometry}

\usepackage[british]{babel}
\usepackage[final]{microtype}

% ===========================
% HEADINGS
% ===========================
\usepackage{titlesec}

% Section: small caps, number in margin
\titleformat{\section}
  {\normalfont\scshape}
  {\llap{\thesection\quad}}
  {0pt}
  {}

% Subsection: small caps sentence case
\titleformat{\subsection}
  {\normalfont\normalsize\scshape}
  {\thesubsection\quad}
  {0pt}
  {}

% Subsubsection: upright
\titleformat{\subsubsection}
  {\normalfont}
  {\thesubsubsection\quad}
  {0pt}
  {}

% Spacing
\titlespacing*{\section}{0pt}{2ex plus 1ex minus .2ex}{1ex plus .2ex}
\titlespacing*{\subsection}{0pt}{1.5ex plus 1ex minus .2ex}{0.5ex plus .2ex}
\titlespacing*{\subsubsection}{0pt}{1ex plus 0.5ex minus .1ex}{0.3ex plus .1ex}

% ===========================
% RUNNING HEADS
% ===========================
\usepackage{fancyhdr}
\pagestyle{fancy}
\fancyhf{}
\fancyhead[LE]{\small\scshape\leftmark}
\fancyhead[RO]{\small\scshape\@title}
\fancyfoot[LE,RO]{\thepage}
\renewcommand{\headrulewidth}{0pt}

% For oneside documents (most papers):
\fancyhead[L]{\small\scshape\leftmark}
\fancyhead[R]{\small\thepage}
\fancyfoot{}

% ===========================
% COLORS & HYPERLINKS
% ===========================
\usepackage{xcolor}
\definecolor{linkmaroon}{RGB}{128, 0, 32}

\usepackage{hyperref}
\hypersetup{
  colorlinks=true,
  linkcolor=linkmaroon,
  citecolor=linkmaroon,
  urlcolor=linkmaroon,
  pdfauthor={Brett Reynolds},
}

\usepackage{orcidlink}

% ===========================
% QUOTATIONS
% ===========================
\usepackage[style=american]{csquotes} % Double quotes outer (Oxford spelling kept via babel)
\MakeOuterQuote{"}

% ===========================
% SEMANTIC MACROS
% ===========================
% Terms (linguistic concepts being introduced/defined)
\newcommand{\term}[1]{\textsc{#1}}

% Mentions (words/expressions under discussion)
\newcommand{\mention}[1]{\textit{#1}}

% Mentions in headings (angle brackets instead of italics)
% Robust: headings are moving arguments (written to .toc/.aux)
\DeclareRobustCommand{\mentionhead}[1]{⟨\textup{#1}⟩}

% Object language (cited forms, foreign words)
\newcommand{\olang}[1]{\textit{#1}}

% Small-caps abbreviations for glosses
\newcommand{\abbr}[1]{\textsc{#1}}

% Cross-linguistic subscript marker
\usepackage{marvosym}
\newcommand{\crossmark}{\textsubscript{\Cross}}

% ===========================
% MATHS AND SYMBOLS
% ===========================
\usepackage{amsmath,amssymb}

% ===========================
% LINGUISTIC EXAMPLES
% ===========================
\usepackage{langsci-gb4e}
\makeatletter
\@ifundefined{noautomath}{}{\noautomath}
\makeatother

% Judgement markers
\newcommand{\ungram}[1]{*\!#1}
\newcommand{\marg}[1]{?\!#1}
\newcommand{\odd}[1]{\#\!#1}

% ===========================
% LISTS
% ===========================
\usepackage{enumitem}
\setlist{nosep, leftmargin=*}
\setlist[enumerate]{label=\arabic*.}
\setlist[itemize]{label=--}

% ===========================
% TABLES & FIGURES
% ===========================
\usepackage{booktabs}
\usepackage{array}
\usepackage{graphicx}
\graphicspath{{figures/}}

% ===========================
% BIBLIOGRAPHY
% ===========================
\usepackage[backend=biber,style=apa,natbib=true,doi=true,isbn=false,url=true]{biblatex}
\IfFileExists{references-standalone.bib}
  {\addbibresource{references-standalone.bib}}
  {\addbibresource{references.bib}\IfFileExists{references-local.bib}{\addbibresource{references-local.bib}}{}}

% Possessive citation: Author's (Year)
\newcommand{\posscite}[1]{\citeauthor{#1}'s (\citeyear{#1})}

% ===========================
% UTILITIES
% ===========================
\usepackage{xspace}
\newcommand{\eg}{e.g.,\xspace}
\newcommand{\ie}{i.e.,\xspace}
\newcommand{\etc}{etc.\xspace}
 in main document
%
% Compiler: XeLaTeX (default) or LuaLaTeX
%           NOT pdfLaTeX - requires fontspec for fonts

% ===========================
% FONTS
% ===========================
\usepackage{fontspec}

\setmainfont{EB Garamond}[
  Numbers=OldStyle,
  Ligatures=TeX,
  BoldFont={EB Garamond},
]

% IPA fallback font
\newfontfamily\ipafont{Charis SIL}
\newcommand{\ipa}[1]{{\ipafont #1}}

% CJK fallback font (Japanese / Chinese / Korean)
\newfontfamily\cjkfont{Hiragino Mincho ProN}
\newcommand{\cjk}[1]{{\cjkfont #1}}

% Lining figures when needed (tables, years in isolation)
\providecommand{\liningnums}[1]{{\addfontfeatures{Numbers=Lining}#1}}

% Monospace
\setmonofont{Inconsolata}[Scale=MatchLowercase]

% ===========================
% PAGE LAYOUT
% ===========================
\usepackage[
  letterpaper,
  inner=1.25in,
  outer=1in,
  top=1in,
  bottom=1.25in,
  marginparwidth=0.6in,
]{geometry}

\usepackage[british]{babel}
\usepackage[final]{microtype}

% ===========================
% HEADINGS
% ===========================
\usepackage{titlesec}

% Section: small caps, number in margin
\titleformat{\section}
  {\normalfont\scshape}
  {\llap{\thesection\quad}}
  {0pt}
  {}

% Subsection: small caps sentence case
\titleformat{\subsection}
  {\normalfont\normalsize\scshape}
  {\thesubsection\quad}
  {0pt}
  {}

% Subsubsection: upright
\titleformat{\subsubsection}
  {\normalfont}
  {\thesubsubsection\quad}
  {0pt}
  {}

% Spacing
\titlespacing*{\section}{0pt}{2ex plus 1ex minus .2ex}{1ex plus .2ex}
\titlespacing*{\subsection}{0pt}{1.5ex plus 1ex minus .2ex}{0.5ex plus .2ex}
\titlespacing*{\subsubsection}{0pt}{1ex plus 0.5ex minus .1ex}{0.3ex plus .1ex}

% ===========================
% RUNNING HEADS
% ===========================
\usepackage{fancyhdr}
\pagestyle{fancy}
\fancyhf{}
\fancyhead[LE]{\small\scshape\leftmark}
\fancyhead[RO]{\small\scshape\@title}
\fancyfoot[LE,RO]{\thepage}
\renewcommand{\headrulewidth}{0pt}

% For oneside documents (most papers):
\fancyhead[L]{\small\scshape\leftmark}
\fancyhead[R]{\small\thepage}
\fancyfoot{}

% ===========================
% COLORS & HYPERLINKS
% ===========================
\usepackage{xcolor}
\definecolor{linkmaroon}{RGB}{128, 0, 32}

\usepackage{hyperref}
\hypersetup{
  colorlinks=true,
  linkcolor=linkmaroon,
  citecolor=linkmaroon,
  urlcolor=linkmaroon,
  pdfauthor={Brett Reynolds},
}

\usepackage{orcidlink}

% ===========================
% QUOTATIONS
% ===========================
\usepackage[style=american]{csquotes} % Double quotes outer (Oxford spelling kept via babel)
\MakeOuterQuote{"}

% ===========================
% SEMANTIC MACROS
% ===========================
% Terms (linguistic concepts being introduced/defined)
\newcommand{\term}[1]{\textsc{#1}}

% Mentions (words/expressions under discussion)
\newcommand{\mention}[1]{\textit{#1}}

% Mentions in headings (angle brackets instead of italics)
% Robust: headings are moving arguments (written to .toc/.aux)
\DeclareRobustCommand{\mentionhead}[1]{⟨\textup{#1}⟩}

% Object language (cited forms, foreign words)
\newcommand{\olang}[1]{\textit{#1}}

% Small-caps abbreviations for glosses
\newcommand{\abbr}[1]{\textsc{#1}}

% Cross-linguistic subscript marker
\usepackage{marvosym}
\newcommand{\crossmark}{\textsubscript{\Cross}}

% ===========================
% MATHS AND SYMBOLS
% ===========================
\usepackage{amsmath,amssymb}

% ===========================
% LINGUISTIC EXAMPLES
% ===========================
\usepackage{langsci-gb4e}
\makeatletter
\@ifundefined{noautomath}{}{\noautomath}
\makeatother

% Judgement markers
\newcommand{\ungram}[1]{*\!#1}
\newcommand{\marg}[1]{?\!#1}
\newcommand{\odd}[1]{\#\!#1}

% ===========================
% LISTS
% ===========================
\usepackage{enumitem}
\setlist{nosep, leftmargin=*}
\setlist[enumerate]{label=\arabic*.}
\setlist[itemize]{label=--}

% ===========================
% TABLES & FIGURES
% ===========================
\usepackage{booktabs}
\usepackage{array}
\usepackage{graphicx}
\graphicspath{{figures/}}

% ===========================
% BIBLIOGRAPHY
% ===========================
\usepackage[backend=biber,style=apa,natbib=true,doi=true,isbn=false,url=true]{biblatex}
\IfFileExists{references-standalone.bib}
  {\addbibresource{references-standalone.bib}}
  {\addbibresource{references.bib}\IfFileExists{references-local.bib}{\addbibresource{references-local.bib}}{}}

% Possessive citation: Author's (Year)
\newcommand{\posscite}[1]{\citeauthor{#1}'s (\citeyear{#1})}

% ===========================
% UTILITIES
% ===========================
\usepackage{xspace}
\newcommand{\eg}{e.g.,\xspace}
\newcommand{\ie}{i.e.,\xspace}
\newcommand{\etc}{etc.\xspace}

\usepackage{tikz}
\usepackage{pgfplots}
\usepackage{longtable}
\setlength{\headheight}{13.6pt}
\pgfplotsset{compat=1.18}
\usetikzlibrary{arrows.meta,shapes,positioning,calc}
\ifdefined\blindsubmission
\hypersetup{
  pdftitle={Definiteness and Deitality in English: A Projectibility-First Account},
  pdfauthor={},
  pdfkeywords={definiteness, article, demonstratives, determinatives, grammatical categories}
}
\else
\hypersetup{
  pdftitle={Definiteness and Deitality in English: A Projectibility-First Account},
  pdfauthor={Brett Reynolds},
  pdfkeywords={definiteness, article, demonstratives, determinatives, grammatical categories}
}
\fi
\urlstyle{same}
\setcounter{biburlnumpenalty}{9000}
\setcounter{biburlucpenalty}{9000}
\setcounter{biburllcpenalty}{9000}
\newcommand{\myorcid}{0000-0003-0073-7195}
\newcommand{\AuthorWithORCID}[2]{#1\,\orcidlink{#2}}

% Macros
\newcommand{\Deital}{\textsc{deital}}
\newcommand{\Nondeital}{\textsc{non-deital}}
\newcommand{\Deitality}{\textsc{deitality}}
\newcommand{\Definiteness}{\textsc{definiteness}}
% Upright (non-italic) slash that still allows a break after it
\DeclareRobustCommand{\upslash}{\/\textup{~\slash~}}
\newcommand{\MetaphysicsClusterCites}{%
  \ifdefined\blindsubmission
    \parencite{slater2015,khalidi2013,khalidi2018nodes,onishi2022,anonymous2026notEveryStableCluster}%
  \else
    \parencite{slater2015,khalidi2013,khalidi2018nodes,onishi2022,reynolds2026notEveryStableCluster}%
  \fi}
\newcommand{\MetaphysicsSelfCite}{%
  \ifdefined\blindsubmission
    \parencite{anonymous2026notEveryStableCluster}%
  \else
    \parencite{reynolds2026notEveryStableCluster}%
  \fi}
\newcommand{\PrepositionCitation}{%
  \ifdefined\blindsubmission
    \parencite{anonymous2025prepositions}%
  \else
    \parencite{Reynolds2025}%
  \fi}


% Title
\title{Definiteness and Deitality in English:\\A Projectibility-First Account}
\ifdefined\blindsubmission
\author{}
\date{}
\else
\author{\AuthorWithORCID{Brett Reynolds}{\myorcid}\thanks{I thank Cris Chatterjee for helpful feedback.}\\[4pt]
Humber Polytechnic \& University of Toronto}
\date{\today}
\fi

\begin{document}

\maketitle

\begin{abstract}
In English grammar, \term{definiteness} has been used for two profiles that are better kept apart. One is semantic: identifiability, uniqueness, anaphoric recoverability, and related discourse conditions. The other is structural, and English grammar has no agreed name for it. I call it \Deitality{}: its prototype is the demonstratives (deictic, and the historical source of \mention{the}), and its core diagnostics involve domain establishment in true partitives, resistance to neutral existential \mention{there} pivots, and ready hosting of identificational material. A narrow bare-\mention{one} pattern supplies corroborating residue rather than a core diagnostic. Definiteness and deitality usually coincide, which is why they get run together, though the alignment is imperfect: weak definites, generic definites, narrative \mention{this}, and proper names each mark a point where they separate.

The account justifies grammatical categories by the inferences they support. Deitality is a stable, historically sourced, learnable grammatical profile. Once a form or construction is located in the profile, the analysis predicts which residues should survive semantic mismatch, which diagnostics should split under prosodic or specificity manipulation, and which variation should appear across dialects and acquisition. The result is a more predictive account of English definiteness. Semantic definiteness and structural deitality are strongly correlated, partly independent, and leaky in principled ways.
\end{abstract}

\noindent\textbf{Keywords:} definiteness; article; demonstratives; determinatives; grammatical categories

\newpage
\section{Introduction}\label{sec:intro}

One inherited label has been covering two profiles. In English grammar, \term{definiteness} names both a semantic pattern of referent identifiability and a structural pattern of distributional behaviour that has no agreed name of its own. The slippage shows up whenever tests for one are read as probes of the other.

The definiteness effect in existential \mention{there} \parencite{milsark1977,McNally2011} is a morphosyntactic test, but it's often used as evidence about uniqueness or familiarity. Semantic tests for identifiability, like the \mention{which?}-question heuristic in \emph{The Cambridge Grammar of the English Language}, ask whether a description identifies the intended referent so that an addressee doesn't have to ask \mention{Which one?} \parencite{HuddlestonPullum2002}; they don't by themselves sort structural profiles.

I call the structural profile \Deitality{}. The new term names an old pattern: \mention{the}, demonstratives, genitive NPs in determiner function, and related constructions pattern together distributionally even when ordinary definiteness semantics loosens. Deitality is characterized by distributional properties that co-occur in English determinatives and in the NP constructions in which they serve as determiners.\footnote{I follow the category--function distinction of \textcite{HuddlestonPullum2002}: \term{determinative} names a lexical category (\mention{the}, \mention{this}, \mention{some}), while \term{determiner} names the function it most often realizes in NP structure. Nothing in the account assumes a determiner-headed nominal.}

The semantic counterpart, \Definiteness{}, is characterized by interpretive properties related to referent identifiability. Deital marking usually tracks semantic definiteness, but the match breaks in both directions: names can be definite without ordinary deital marking, and weak definites can be deital without ordinary token-level definiteness. The correlation is strong, but the two aren't identical. Nor is deitality just the old strong/weak contrast, specificity, familiarity, or uniqueness under a new name; those notions target semantic, discourse, or construction-specific regions of the pattern.

\textcite{lyons1999} is the closest predecessor. He treats definiteness as a grammatical category tied to the grammaticalization of identifiability, with its own syntax, distribution, diachrony, and cross-linguistic profile. I accept that separation between grammatical marking and the pragmatic concept it grammaticalizes. The claim here is narrower, but it adds a different kind of test: English contains a deital profile within the determination system, anchored in demonstrative-derived and determiner-function forms, whose ranked diagnostics support predictions about partitives, existential pivots, hosting, mismatch residue, dialect variation, and acquisition.

Proper names show the split in a concrete way. Bare singular names such as \mention{John} are semantically definite but lack ordinary deital marking, while plural family names such as \mention{the Smiths} require it. So-called \enquote{weak definites} like \mention{go to the hospital}\footnote{I retain the standard label; §\ref{sec:cases} argues that these constructions use deital form while lacking ordinary token-level definiteness.} combine definite morphology with semantically indefinite reference. \enquote{Generic definites} like \mention{The tiger is endangered} use definite singular form for kind reference rather than individual-level definiteness. These form--function mismatches are calibration points \parencite{FrancisMichaelis2003}: places where the semantic and structural profiles detach just enough to expose the structural pattern beneath the inherited label.

The analysis uses one idea from natural-kind theory: a category is projectible when observed cases license expectations about unobserved ones \parencite{goodman1983}. Section~\ref{sec:hpc} develops the Boydian and post-HPC discipline behind that claim; the point here is that deitality has to support predictions across constructions, semantic mismatches, manipulations, and borderline cases \parencite{boyd1991,boyd1999}.

The sources and stabilizers that matter here operate at distinct timescales, from demonstrative grammaticalization to acquisition and usage. Because deitality and definiteness are supported by partly different pressures, the overlap between them should hold in some regions and loosen in others.

The payoff is empirical first. The projectibility-first framework explains why deitality and definiteness align without being identical, why weak and generic definites exist, and why proper names pattern as they do. It also tells us what to expect once we locate a determinative, noun phrase, or construction in the deital profile: which residues should survive semantic mismatch, which diagnostics should split under manipulation, and which parts of the system should vary across prosodic, dialectal, and developmental conditions. Theoretically, it offers a principled way to handle internal gradedness without making the category unreal: how central an item is, how strong the evidence is, and where the boundary falls are distinct questions.

The paper proceeds as follows. Section~\ref{sec:hpc} introduces the projectibility-first framework and justifies its application to grammatical categories. Section~\ref{sec:semantic} characterizes definiteness as a semantic profile. Section~\ref{sec:diagnostics} develops the diagnostics for deitality as morphosyntactic tests. Section~\ref{sec:mechanisms} identifies sources and stabilizers of the deital profile and derives empirical predictions. Section~\ref{sec:cases} applies the framework to weak and generic definites and narrative \mention{this}. Section~\ref{sec:objections} addresses edge cases and objections. Section~\ref{sec:conclusion} concludes.

\section{Projectible profiles and grammatical categories}\label{sec:hpc}

The proposal to treat deitality as a projectible grammatical profile requires justification. The framework is Boydian in motivation: it rejects necessary-and-sufficient definitions and asks what categories let us infer. But it also incorporates a later discipline from natural-kind theory: stability, network order, maintenance, and strict homeostasis aren't the same achievement.

\subsection{From essence to projectibility}

The motivating contrast comes from the species problem. Essentialist definitions of species sought a set of necessary and sufficient conditions for membership: morphological similarity, ability to interbreed, shared ancestry, or some other privileged property. Each candidate failed. Ring species can interbreed with adjacent populations but not with distant ones. Asexually reproducing organisms don't interbreed at all. Morphologically similar organisms can be reproductively isolated. The natural world resisted any single fixed cutoff \parencite{mayr1942}.

\textcite{boyd1991,boyd1999} rejects essentialism. A species is characterized by a cluster of properties (morphological, genetic, behavioural) that co-occur more often than chance would predict, with no single property or fixed list required for membership. What matters for scientific practice is projectibility: for \textcite{goodman1983}, a category supports projections from observed cases to unobserved ones. In Boyd's HPC framework, that means observing some parts of the profile licenses expectations about others. The projectibility is also field-relative: a category earns its place by supporting the inferences a practice needs. The category is useful because it fits inferential practice to real causal structure, not because it mirrors an essence.

This gives three commitments that are useful for grammar. First, the category has a contingent property profile: a family of properties that statistically co-occur. Second, the profile has enough order or maintenance to keep the co-occurrence from being accidental. Third, the profile is projectible for some analytical purpose. Because these conditions aren't essentialist, the resulting category can have internally graded structure and prototype effects without collapsing into arbitrariness. That doesn't have to make the boundary itself fuzzy. It means that our evidence for locating the boundary can be uneven, and that category members can be more or less central once they're inside it.

\subsection{Separating stability, maintenance, and control}

The word \term{homeostatic} shouldn't do more work than the evidence warrants. In the HPC-derived debate, a profile can be stable without a known maintainer; it can be network-ordered without feedback; and it can be maintained without being a corrective-control system \parencite{slater2015,khalidi2013,khalidi2018nodes,onishi2022}. A stricter diagnostic asks five questions in order \MetaphysicsSelfCite{}:

\begin{enumerate}[label=(\roman*)]
\item What's the projection target?
\item What stable profile supports that projection?
\item What ordering relations explain why the projection isn't accidental?
\item What maintains the profile across relevant timescales?
\item Is the maintainer a corrective-control process that preserves a higher-scale relation through lower-scale variation?
\end{enumerate}

This paper doesn't need the fifth step. The central linguistic claim is that deitality is a projectible grammatical profile, not that English determinatives form a homeostatic control system. Its projection targets include more than the diagnostics that first identify it. They include external and manipulable outcomes: which residues survive when semantic definiteness is relaxed, which contrasts are rescued by prosody, which frames vary across dialects, and which restrictions should be acquired earlier than others. A category is useful only if it licenses those further inferences.

\subsection{Why grammatical categories need this framework}

The parallel to linguistics is direct. Just as biologists sought an essence for species, linguists have sought one for definiteness. Decades of research have tried to define \enquote{definite noun phrase} by semantic essence, usually uniqueness \parencite{russell1905,hawkins1978} or familiarity \parencite{christophersen1939,heim1982}. But counterexamples proliferate. Weak definites like \mention{listen to the radio} use \mention{the} without uniqueness \parencite{carlsonsussman2005,aguilarguevarazwarts2010}. Generic definites like \mention{The lion is noble} refer to kinds, not individuals \parencite{carlson1977,krifka2004}. Cross-linguistically, definite markers appear in contexts where their semantic contribution seems redundant or absent \parencite{lyons1999}.

Far from being peripheral exceptions, these systematic patterns suggest the category lacks a simple essence. The attempt to define a morphosyntactic object, the category of noun phrases marked with definite morphology, using purely semantic criteria has produced a category that doesn't align with the empirical facts.

The projectibility-first framework offers a way forward. We can characterize deitality and definiteness as distinct profiles and ask what each lets us infer. The strong correlation between them is real and grounded in grammar, discourse, history, and acquisition, but the profiles can come apart in systematic ways \parencite{FrancisMichaelis2003}.

Cross-linguistic work on two kinds of definites gives an important comparison point. \textcite{schwarz2009,schwarz2013} argues that German weak and strong article forms track different semantic relations, roughly uniqueness and anaphoric familiarity; related work extends the comparison to languages without ordinary articles and to languages where the analysis is contested \parencite{jenks2018,arkohMatthewson2013,bombi2018}. This literature is an ally because it shows that the definiteness domain can split grammatically. The account here asks a different question. It doesn't replace the uniqueness--familiarity contrast; it asks whether English forms and constructions participate in a ranked distributional profile whose diagnostics split under prosody, specificity, semantic mismatch, dialectal variation, and acquisition. On this account, the Schwarz-style weak/strong contrast falls within deital territory rather than exhausting it: the question is how semantically different definite markers pattern on structural diagnostics.

This move has precedent. Prototype theory in cognitive linguistics has long recognized that categories have graded structure \parencite{rosch1978,lakoff1987}. But prototype theory often stops at similarity structure. Boyd's HPC framework adds a causal account of why certain properties cluster and why that clustering supports reliable inferences beyond the data that first identified it.

The framework has already been applied to linguistic phenomena. \textcite{miller2021} argues that word-kinds like \mention{table} and \mention{Paris} are HPCs, defined by clusters of phonological, orthographic, and semantic properties maintained by cognitive and social mechanisms. This paper extends that style of analysis from lexical ontology to grammatical categories, but with a stricter vocabulary. Where Miller asks what makes tokens instances of the same word-kind, I ask what distinguishes deital from non-deital determinatives: why \mention{the}, \mention{this}, and genitive NPs in determiner function pattern together distributionally while \mention{a}, \mention{some}, and \mention{many} form a contrasting profile.

\section{Definiteness as a semantic profile}\label{sec:semantic}

Before characterizing deitality, I need to clarify what definiteness is, semantically and pragmatically, so the contrast with deitality is clear. The literature presents two major traditions: uniqueness theories and familiarity theories. Nearby work also separates semantic properties often grouped under definiteness, including definiteness and determinacy \parencite{coppockBeaver2015}. Rather than adjudicating between these approaches, I argue that they capture important properties in the definiteness profile.

Uniqueness theories originate with \textcite{russell1905}. A definite description \mention{the F} presupposes that exactly one entity satisfies \mention{F} in the contextually relevant domain. Later work refined this to allow pragmatic restriction: uniqueness holds relative to a \enquote{P-set} or shared frame of reference \parencite{hawkins1978}. This approach handles \enquote{larger situation} uses like \mention{the sun} and anaphoric uses where prior discourse has narrowed the context to a single salient entity.

Familiarity theories stem from \textcite{christophersen1939} and were formalized by \textcite{heim1982}. Definite noun phrases are subject to a familiarity condition: they have to refer to an entity already established in the discourse model. Indefinites introduce new referents. This framework excels at explaining discourse anaphora, where a definite's felicity depends directly on having a linguistic antecedent.

Both theories have problems. Familiarity isn't necessary: \mention{the first person on Mars} can be used felicitously even when the referent is discourse-new. Uniqueness isn't sufficient: in contexts with two familiar but non-unique entities, neither \mention{the N} works smoothly. And weak definites challenge both: \mention{go to the hospital} uses \mention{the} when the hospital is neither uniquely specified nor discourse-familiar \parencite{Abbott2004,Birner2013}.

I propose that uniqueness and familiarity are both properties in the definiteness profile, along with anaphoric recoverability and, more broadly, identifiability. These properties co-tend because of discourse-pragmatic pressures: common ground management makes unique entities easier to track; anaphoric devices rely on familiarity; topics bias toward identifiable referents. That's why uniqueness and familiarity theories map different regions of one semantic profile rather than competing essences. The properties overlap heavily (a discourse-old, unique referent satisfies both), but they dissociate in predictable ways. Weak definites have identifiability without token-level uniqueness. Cataphoric definites have uniqueness without prior familiarity. The semantic notion of definiteness is profile-like rather than essence-like.

Definiteness is a profile centred on interpretation: how referents are mentally represented and tracked in discourse. The core properties are identifiability, uniqueness, and anaphoric recoverability. Definite referents are often marked deitally, but that association is an empirical correlation to be explained, not a defining property of the semantic profile. The next sections establish the deital profile through convergent diagnostics (§\ref{sec:diagnostics}) and identify the distinct processes maintaining it (§\ref{sec:mechanisms}). The imperfect alignment between profiles, cases where definiteness and deitality dissociate, is the focus of §\ref{sec:cases}.

\section{Diagnostics for deitality}\label{sec:diagnostics}

I characterize deitality as a structural profile whose core comprises distributional restrictions. It also shows a strong empirical tendency to mark semantically definite referents, but that tendency belongs to the correlation between profiles rather than to the definition of deitality itself. The profile is level-sensitive. Sometimes the relevant object is a determinative lexeme (\mention{the}, \mention{this}); sometimes it's a genitive NP in determiner function; sometimes it's an NP in a particular construction; sometimes it's a diachronic lineage whose forms preserve inherited distributional residue. Deitality is a category distributed across levels rather than a single feature borne by one grammatical object in all contexts. It's the category that answers a syntactician's question: which forms and constructions support the same projectible distributional inferences?

This section develops three uneven diagnostics plus one residue pattern: existential \mention{there}, domain-establishing \mention{Y} complements in true partitives, and identificational hosting. The partitive diagnostic is the strongest syntactic test. Existential \mention{there} depends on construction and information structure. Hosting is mixed: it reflects morphosyntax, specificity, and discourse anchoring. Bare-\mention{one} behaviour is fragile corroborating residue, not a core diagnostic. The tests are useful precisely because they don't reduce to one another.

\subsection{Existential \mentionhead{there}}

The definiteness effect in existential \mention{there} constructions is well documented \parencite{milsark1977,McNally2011}. Under neutral prosody, deital determinatives resist the pivot position:

\ea 
\label{ex:there-neutral}
    \ea \textit{\textsuperscript{*}There's the key on the table.} (neutral existential reading)
    \ex \textit{\textsuperscript{*}There are those books I mentioned.}
    \ex \textit{There's a key on the table.} (non-deital, acceptable)
    \z
\z

\noindent The effect is sensitive to prosody: list/presentational intonation can license otherwise dispreferred pivots \parencite{BirnerWard1998}:

\ea 
\label{ex:there-list}
    \ea \textit{There's THAT book again.}
    \ex \textit{Well, there's THE KEYS, the wallet...}
    \z
\z

\noindent The initial \mention{there} in such cases can be heard as locative rather than expletive, which makes the intended presentational reading harder to obtain.%
\ifdefined\blindsubmission
\footnote{Thanks to a reader for pressing this point.}
\else
\footnote{Thanks to Cris Chatterjee for pressing this point.}
\fi

I treat this as a controllable neutralization: under neutral prosody, deital determinatives resist the pivot position; marked prosody can ameliorate the effect. I calibrate the diagnostic to neutral prosody.

This calibration matters. Existential \mention{there} has semantic and information-structure inputs, so list readings, recurrence readings, locative/expletive ambiguity, and narrative frames can affect judgments. I use it as a ranked cue under neutral prosody and non-narrative frames, not as sufficient evidence for deitality.

\textcite{McNally2011} is especially important because the definiteness restriction is treated there as more than a simple surface-form ban. McNally reviews semantic accounts on which the pivot has to be property-like and pragmatic accounts on which the pivot is constrained by discourse status or information structure; her own earlier analysis is explicitly mixed. I accept that lesson. The \mention{there} diagnostic doesn't by itself show a morphosyntactic category distinct from definiteness. Its value here is comparative: it patterns with the profile under neutral conditions, but it dissociates from the harder partitive diagnostic under prosodic manipulation. That contrast makes the cue useful, but it isn't decisive on its own; §\ref{sec:objections} returns to the semantic-bundle worry and adds acquisition and weak-definite tests.

\subsection{Partitive complements}

In true partitive constructions (\mention{X of Y} with subset semantics), the relevant diagnostic belongs to the Partitive Constraint tradition \parencite{ladusaw1982,abbott1996,Barker1998,Hoeksema1996}. The issue is the availability of a delimited domain rather than \mention{the}-marked definiteness as such. I state the diagnostic conservatively: the \mention{Y} complement has to be definite-like or independently domain-establishing.

\ea 
\label{ex:part-good}
    \ea \textit{Two\upslash several\upslash many of the students left.}
    \ex \textit{Two\upslash several\upslash many of these\upslash my students left.}
    \z
\z

\ea 
\label{ex:part-bad}
*\textit{Several of some students left.}
\z

\noindent This selectional restriction is stable and isn't ameliorated by prosody. The partitive literature contains complications, including examples such as \mention{one of several options}, \mention{one of many problems}, and \mention{two of a number of applicants} \parencite{ladusaw1982,abbott1996}. Those cases show that the domain-establishing condition can be satisfied without \mention{the}. They also show why the diagnostic shouldn't be stated as a simple article-definiteness test. The narrower contrast is still useful: ordinary weak indefinites like \mention{some students} don't behave like \mention{the students}, \mention{these students}, or genitive NPs as partitive complements. Attempts to use weak \mention{Y} are either rejected outright, receive a specialized set-denoting interpretation, or are reanalyzed as measure/fraction constructions (\mention{half of a cake}), which have different semantics.

The \mention{Y} slot is what matters. Elements in \mention{X}, including quantifiers like \mention{all} or \mention{both} when they appear as modifiers, are irrelevant to this diagnostic.

\subsection{Identificational hosting}

I use \term{identificational hosting} as a cover term for environments that need a discourse-salient, identifiable NP as host: nonrestrictive modification, \mention{as for} topics, and specificational subjects (the copular type of \mention{The culprit is John}, distinct from the identificational \mention{That is John}). The requirement on the host NP is semantic and information-structural; the deital diagnostic is distributional. NPs with deital determiners satisfy these environments most readily, while NPs with weak determiners need independent anchoring, especially in topic and specificational frames. Nonrestrictive relatives are more permissive when the indefinite is independently anchored \parencite{Lambrecht1994}:

\ea 
\label{ex:host-good}
    \ea \textit{The\upslash this\upslash John's book, which I bought yesterday, is excellent.}
    \ex \textit{As for the\upslash this\upslash John's report, I filed it.}
    \ex \textit{The culprit is John.}
    \z
\z

\ea 
\label{ex:host-mixed}
    \ea \textit{A student of mine, who you met, just won a prize.} (anchored indefinite)
    \ex \textsuperscript{??}\textit{As for a report, I filed it.}
    \ex \textsuperscript{??}\textit{A culprit is John.}
    \z
\z

\noindent The degradation correlates with information structure. These host constructions presuppose that the referent is already salient or identifiable to the hearer. Weak indefinites introduce new referents and clash with this presupposition unless other material supplies anchoring. This makes the diagnostic mixed rather than purely structural: examples such as \mention{A colleague of mine, who you may remember from York, emailed me} improve because specificity and discourse anchoring supply part of what the construction requires. The prediction is graded, not categorical.

\subsection{Narrow \mentionhead{one}-substitution residue}

Non-deital determinatives freely allow bare anaphoric \mention{one} across clauses:

\ea 
\label{ex:one-weak}
\textit{I bought a red pen and you bought one too.}
\z

\noindent Deital antecedents resist bare \mention{one} on readings where the continuation would have to preserve the already-determined token rather than introduce another instance of the nominal property:

\ea 
\label{ex:one-strong}
\textsuperscript{*}\textit{I bought the\upslash this\upslash my red pen and you bought one too.} (intended: the same contextually determined pen)
\z

\noindent Determiner function + \mention{one} (\mention{this one}, \mention{the one}) is grammatical but represents a different construction: \mention{one} is a noun with its own determiner, not a determinative phrase in fused-determiner--head function (one word filling the determiner and head positions together, as in \mention{I'll take these}) \parencite[332--333]{HuddlestonPullum2002}. I exclude this construction from the residue pattern.

The limitation is built into the analysis. In \CGEL{} terms, bare \mention{one} is a nominal anaphor; it isn't a device for recovering a whole determiner--head package. Many speakers accept examples like \mention{I found the perfect apartment and she found one too}, where \mention{one} recovers the nominal property rather than the determined phrase. The test is useful only when the context makes the determined-token reading salient and the property-anaphoric reading pragmatically unavailable. I treat it as a residue to discuss rather than as an inventory column.

\subsection{Summary: Converging evidence for a grammatical profile}

These diagnostics converge on a morphosyntactic profile, but they aren't equally stable. Prototypical deital determinatives resist existential pivots, satisfy the domain-establishing \mention{Y} condition in partitives, and host identificational constructions naturally. They also show the narrow \mention{one}-substitution residue just described. Prototypical non-deital determinatives show the opposite pattern on the stronger diagnostics. No single diagnostic defines the profile. Partitive \mention{Y}-licensing has the greatest syntactic weight; existential \mention{there} and hosting show how the profile interacts with information structure; \mention{one}-substitution is too fragile for category placement on its own. What matters is ranked convergence across tests, with residue patterns used only as corroboration.

In applying the tests (see Appendix~\ref{sec:inventory}), many determinatives show partial profile convergence. Items like \mention{each}, \mention{every}, \mention{all}, and \mention{both} resist existential pivots (patterning deitally) but fail other diagnostics (patterning non-deitally). The projectibility-first account predicts this internal gradedness: different degrees of property convergence. The diagnostic convergence maps a space where \mention{the} and demonstratives occupy the prototypical centre while other items supply less decisive evidence about where the boundary falls. I use \term{Mixed} in the inventory as an evidential status: it marks partial convergence and uncertainty about placement, not literal half-membership.

This property pluralism is central to the framework: just as biological species are characterized by multiple co-tending traits (morphology, genetics, behaviour) with none metaphysically prior, deitality is characterized by multiple co-tending distributional properties with none foundational. The diagnostics differ, and their dissociations are informative rather than embarrassing.

These diagnostics target distribution, not semantic contribution. Deital forms often mark identifiable referents, but that correlation is an empirical fact to explain rather than a defining condition. A determinative can be deital while contributing non-prototypical semantics: weak definites lack token-level uniqueness; generic definites refer to kinds. The morphosyntactic profile is what remains stable across this semantic variation.

Section~\ref{sec:mechanisms} identifies the processes that source and stabilize this diagnostic profile.

\section{Sources, stabilizers, and predictions of the deital profile}\label{sec:mechanisms}

The central claim of the projectibility-first account is that property convergence isn't accidental. For deitality, the evidence is strongest for a diachronic source and a learning filter: grammaticalization creates much of the profile, and acquisition plausibly helps reproduce it in each generation. Other processes (interactive alignment, prestige selection, and iterated transmission) are plausible stabilizers of particular parts of the profile, especially at gradient margins. I treat them as ranked contributors rather than as proof of a single rich homeostatic mechanism. For definiteness, the relevant pressures are discourse-pragmatic: information structure and common ground management create correlations among identifiability, uniqueness, and anaphoric recoverability.

The previous section established diagnostic properties that converge in the deital profile. This section identifies sources and stabilizers operating at distinct timescales and derives falsifiable predictions about synchronic grammar.

\subsection{Sources and stabilizers across timescales}\label{subsec:five-mechanisms}

Grammaticalization creates much of the profile over phylogenetic time. Cross-linguistically, definite markers arise primarily from demonstratives \parencite{greenberg1978,diessel1999}. As a demonstrative grammaticalizes (weakening phonologically and broadening functionally), some distributional properties can persist while the speaker-anchored deictic meaning bleaches. The conjecture is selective, not automatic: properties anchored in constructional selection or entrenched morphosyntactic frames should be better candidates for persistence than properties that depend on live speaker-centred deixis. Reanalysis happens utterance by utterance in contexts where the old demonstrative still satisfies the local frame. As the semantic function generalizes from pure deixis to broader identifiability marking, the frame can remain. I call this \term{property dragging}: grammaticalization can preserve formal behaviour while semantic function shifts.

English \mention{the} shows this pattern. Historically derived from a demonstrative, it resists existential pivots, appears as a domain-establishing \mention{Y} in partitives, and hosts identificational constructions, all properties compatible with its deictic source. Weak definites, where \mention{the} appears without uniqueness, represent cases where the morphosyntactic frame has outlived the original semantic restriction. The grammaticalization account predicts that items sharing a deictic source should share distributional properties most reliably where those properties are tied to constructional frames. English offers a modest lineage trace of this kind: the /ð/-initial determinatives (\mention{the}, \mention{this}, \mention{that}, \mention{these}, \mention{those}) are all deital and all belong to the inherited Germanic demonstrative family. I return to the limited evidential role of this trace in §\ref{sec:objections}.

First language acquisition filters the hypothesis space in childhood. Each new generation has to acquire the grammatical system, and cultural transmission favours patterns that are stable and learnable \parencite{kirby2014}. Children acquire deital determinatives early, often overgeneralizing \mention{the} to contexts where adults use \mention{a} \parencite{maratsos1976,rozendaalbaker2008}. This doesn't have to be pure semantic confusion. It's compatible with a structural-learning story: children identify a high-frequency form (\mention{the}) plus parts of its distributional frame and apply it broadly before refining the semantic conditions.

The stronger developmental claim remains conditional. Children may master some positions for \mention{the} before the full pragmatic conditions on adult use. The overgeneralization asymmetry is still positive evidence for a structural-learning story: children seem to latch onto \mention{the} and parts of its frame before the adult semantic conditions are fully in place. But the relative timing of specific restrictions, including partitive \mention{Y}-licensing, still needs direct evidence. If borne out, this would support a bootstrapping path in which distributional frames help learners infer semantic function, and vice versa. Until that evidence exists, bootstrapping is a forecast, not a result. Each generation re-creates the deitality category from distributional and interpretive evidence, so children shouldn't learn pivot resistance, partitive \mention{Y}-licensing, and identificational hosting as wholly unrelated facts.

Other familiar stabilizers are more programmatic here. Interactive alignment, prestige-weighted selection, and iterated cultural transmission are real mechanisms in language use and change \parencite{PickeringGarrod2004,labov2001,eckert2000,kirby2014}, but they're too general to explain deitality by themselves. To make them deitality-specific, we'd need signatures such as repair sequences targeting deital misuse, socially stratified article or demonstrative variants tied to the diagnostics, or transmission effects that preserve selectional restrictions more strongly than local weak-definite frames. The evidence reviewed below uses these processes mainly to locate where variation should appear. Grammaticalization and acquisition carry the core maintenance claim.

\subsection{Causal architecture: How the sources and stabilizers interact}


Figure~\ref{fig:deital-maintainers} schematizes where stabilization could enter the deital profile. Grammaticalization creates much of the initial bundle; the other processes are candidate sites where the profile can be reproduced, filtered, or reshaped at different timescales.

\begin{figure}[htbp]
\centering
\begin{tikzpicture}[
  font=\small,
  node distance = 10mm and 18mm,
  origin/.style = {rectangle, draw, rounded corners,
                   minimum width=40mm, align=center, text depth=0.5ex,
                   fill=blue!20, minimum height=12mm},
  mech/.style = {rectangle, draw, rounded corners,
                 minimum width=38mm, align=center, text depth=0.5ex,
                 fill=blue!10, minimum height=16mm},
  cluster/.style = {rectangle, draw, rounded corners, thick,
                    minimum width=42mm, minimum height=32mm, align=center},
  prim/.style = {->, >=latex, thick},
  sec/.style = {->, dashed, >={stealth}}
]

% Top: origin process
\node[origin] (gram) at (0,0) {Grammaticalization\\(phylogenetic, centuries)};

% Middle: Property profile being maintained
\node[cluster, below=20mm of gram] (props) {
  \textbf{Deital profile}\\[2mm]
  \small
  partitive \mention{Y}-licensing\\
  \mention{There}-resistance\\
  Identificational hosting\\
  narrow \mention{one} residue
};

% Left: Maintenance processes
\node[mech, left=of props, yshift=12mm] (iter) {Iterated\\transmission\\(generations)};
\node[mech, left=of props, yshift=-12mm] (prestige) {Prestige\\selection\\(decades)};

% Right: Maintenance processes
\node[mech, right=of props, yshift=12mm] (align) {Interactive\\alignment\\(seconds--minutes)};
\node[mech, right=of props, yshift=-12mm] (acq) {First language\\acquisition\\(childhood)};

% Grammaticalization creates the bundle
\draw[prim] (gram.south) to node[right, font=\scriptsize] {creates} (props.north);

% Four maintenance processes (all simultaneous)
\draw[prim] (iter.east) to (props.west |- iter.east);
\draw[prim] (prestige.east) to (props.west |- prestige.east);
\draw[prim] (align.west) to (props.east |- align.west);
\draw[prim] (acq.west) to (props.east |- acq.west);

% Feedback: stable profile enables acquisition
\draw[sec, rounded corners] (props.south) to [out=270, in=270] 
  node[below, font=\scriptsize] {stabilizes} (acq.south);

\end{tikzpicture}
\caption{Candidate sources and stabilizers of the deital profile. Grammaticalization (top) created much of the initial property bundle through demonstrative-to-definite-marker pathways. The other processes identify possible sites of reproduction, filtering, or reshaping: iterated transmission across generations; prestige-weighted selection over decades; interactive alignment in real-time conversation; first language acquisition in childhood. The dashed arrow marks feedback between learnability and stability.}
\label{fig:deital-maintainers}
\end{figure}

The feedback arrow isn't an explanation by itself. Stable things get learned, and learned things can stay stable. The deital profile needs a more specific persistence story: high token frequency in the article/demonstrative family plus hard selectional anchors such as partitive \mention{Y}-licensing make it a better candidate for reproduction than lower-frequency or socially fragile patterns.

\subsection{Why deitality and definiteness correlate imperfectly}\label{subsec:autonomy}

Grammaticalization and acquisition explain why deitality exists as a stable morphosyntactic profile. But why does it correlate with definiteness, and why isn't the correlation perfect? The correlation arises because the deictic sources of definite markers were themselves sensitive to information structure. Demonstratives mark salient, often identifiable referents. As they grammaticalize, this functional tendency becomes conventionalized. Speakers and hearers coordinate on using \mention{the} for identifiable referents because that's what the form conventionally signals.

This coordination is maintained by communicative efficiency. Speakers remember frequent pairings, \mention{the} plus unique/familiar referents, and hearers expect them. When speakers violate expectations, hearers accommodate by adjusting their domain of evaluation, coercing to specific readings, or reanalyzing the structure. This creates a stable equilibrium where form and function align statistically.

But the correlation remains imperfect because grammaticalization preserves formal residue. Weak definites like \mention{take the bus} show the pattern especially clearly: a conventionalized \mention{the}-frame can persist while ordinary token-level definiteness weakens. Section~\ref{sec:cases} develops the case in detail.

The mismatch also reflects constructional autonomy. Once a morphosyntactic pattern is established, it can be recruited for new functions \parencite{FrancisMichaelis2003}. \mention{The poor} and \mention{the unknown} as fused modifier--head NPs use \mention{the} to signal definiteness at the constructional level, not the lexical level. The frame itself becomes meaningful, semi-independent of the core function of \mention{the}.

\subsection{Predictions from the profile}

A category earns its place by supporting inductive and explanatory practice. Deitality earns that place when partial evidence about its profile lets us predict untested behaviour. The predictions below make that payoff explicit: they target prosody, specificity, dialectal interpretation, and acquisition.

The predictions have different roles. Prosody and specificity are manipulation tests: they ask which diagnostics split when information structure or discourse anchoring is changed. The dialectal prediction asks where pathway-conditioned variation should appear. The acquisition prediction asks whether children learn hard selectional restrictions earlier than softer discourse-pragmatic ones. The force is cumulative. The prosody split is a within-account falsifier, but it doesn't by itself defeat every semantic-bundle rival; §\ref{sec:objections} returns to that point and adds acquisition and weak-definite existential tests.

\subsection{Prosodic manipulation dissociation}\label{sec:prosody}

Prosody should rescue existential pivots, but not partitive \mention{Y} complements. That contrast isn't inferred from prosody alone. Partitive \mention{Y}-licensing is an internal complement relation in an \mention{X of Y} construction; it appears in written examples and embedded clauses, where intonation supplies no rescue path. Existential \mention{there}, by contrast, is separately sensitive to presentation, list readings, recurrence readings, and locative/expletive ambiguity. Prosody can rescue what discourse mediates; prosody can't rescue what selection demands.

The relevant test cases are:
\ea
    \ea \textit{\textsuperscript{*}There's the key on the table.} (neutral existential reading: unacceptable)
    \ex \textit{Well, there's THE KEY, the wallet...} (list prosody: acceptable)
    \ex \textit{\textsuperscript{*}Several of some students passed.} (neutral prosody: unacceptable)
    \ex \textit{\textsuperscript{*}SEVERAL of SOME students passed.} (contrastive prosody: still unacceptable)
    \z
\z

\noindent If contrastive prosody rescues partitive \mention{Y} complements to the same degree it rescues existential pivots, the projectibility-first account fails at this point. If the split holds, it supports the proposed causal contrast between discourse-pragmatic and syntactic constraints. It doesn't by itself settle the broader rivalry with semantic-bundle accounts; that stronger discriminating work comes from the cumulative pattern.

\subsection{Specificity forcing partial amelioration}\label{sec:specific}

Forcing specificity should improve indefinites in identificational environments without fully licensing them. A semantically specific indefinite keeps its non-deital morphosyntactic profile. The improvement should be measurable but fall short of canonical deital hosts because the morphosyntactic mismatch persists.

The relevant test cases are:
\ea
    \ea \textit{\textsuperscript{??}As for a report, I filed it.}
    \ex \textit{\textsuperscript{?}As for a certain report, I filed it.} (improved but still marked)
    \ex \textit{As for the\upslash this report, I filed it.} (fully acceptable baseline)
    \z
\z

\noindent Acceptability ratings should show this ordering: baseline indefinite < specific indefinite < deital baseline. If specific indefinites reach deital baseline acceptability, the morphosyntactic/semantic distinction collapses. If they show no improvement, semantic contribution is irrelevant to hosting. Partial improvement confirms that both morphosyntactic form and semantic contribution matter.

\subsection{Dialectal frame interpretation}

The broad dialectal prediction is layered. The clean discriminator is specified in advance: a variety that allowed weak indefinites as ordinary partitive \mention{Y} complements would show a deeper grammar difference than a variety that differs only in institutional frames such as \mention{in hospital}. Other variation may affect article realization or presence, reinforced demonstratives, and demonstrative paradigms \parencite{Jones2002,RuppTagliamonte2019ZeroArticle,RuppTagliamonte2019ThisHere,HazenHamiltonVacovsky2011}. These are hypotheses about where to look, not an inference from stability to depth and back again.

The narrower standard-variety test concerns institutional frames. Dialectal variation in those frames should preserve the institutional/activity reading in the licensed frame, while demonstratives and genitive NPs should tend to force ordinary referential readings. The relevant test isn't whether demonstratives are accepted. Demonstratives are accepted: \mention{this hospital} and \mention{that hospital} are ordinary referential NPs. The question is whether the weak-definite or institutional reading survives. If the standard-variety difference is tied to conventionalized frames rather than free alternation in determinative syntax, replacing the frame's licensed form with a demonstrative should usually shift the interpretation from `receiving institutional care' to `located in a particular hospital'.

The relevant test cases are:
\ea
    \ea British English: \textit{She's in hospital.} (institutional/activity reading)
    \ex American English: \textit{She's in the hospital.} (institutional/activity reading available)
    \ex Both dialects: \textit{She's in this hospital.} (acceptable, but preferentially ordinary locative/referential)
    \z
\z

\noindent If demonstratives preserve the institutional/activity reading just as well as the dialectally licensed weak-definite frame, then the account has overstated the special status of the frame. If they instead force a referential-location reading while the licensed frame preserves the institutional reading, the profile supports the intended inference: dialects vary over frames, not freely over deital syntax.

\subsection{Acquisitional asymmetries}

Children should show developmental asymmetries reflecting the sources and stabilizers at work. If the profile is maintained by acquisition, its internal structure should be visible in the learning path. Children should master hard syntactic constraints early because these lack prosodic repair paths. Softer discourse-pragmatic constraints should be mastered later.

The relevant test cases are:
\ea
    \ea Children should correctly reject ungrammatical partitives like \mention{*several of some students} earlier than they master adult-like neutral-prosody restrictions on existential pivots.
    \ex Children may over-accept deitals in existential \mention{there} contexts, especially with list or recurrence intonation, before mastering the adult-like neutral-prosody constraint (a softer, discourse-pragmatic rule).
    \z
\z

\noindent No particular age estimate is offered. The prediction is an ordering: partitive \mention{Y} violations should be rejected before existential violations under neutral prosody are handled in fully adult-like ways. The gap reflects the availability of prosodic rescue for existentials but not partitives. The experiment is demanding because partitive syntax is late enough to require targeted testing, and partitive uses are sparse in spontaneous child and child-directed corpora \parencite{Stickney2007,BloomWynn1997}. The claim needs targeted comprehension work rather than casual production counts.

\subsection{Distinguishing the projectibility-first account from alternatives}

These targeted predictions also separate the account from nearby alternatives. Milsark's strong/weak account captures the existential contrast, but doesn't predict the split among partitive \mention{Y}-licensing, hosting, and narrative \mention{this}. A uniqueness/familiarity account expects prosodic rescue to track referent identifiability more evenly across constructions, which leaves it with trouble explaining why list prosody improves \mention{there} pivots but not non-domain-establishing partitive \mention{Y} complements. A specificity or D-linking account predicts improvement under specificity forcing, but doesn't predict the persistent gap between specific indefinites and deital baselines in identificational hosting.

The more local alternatives miss different projections. A frame-by-frame account can describe each local frame, but it doesn't by itself predict projection across diagnostics: knowing that an item is deital should support expectations about partitives, hosting, prosody, and mismatch residue. A purely diachronic account explains why the family exists historically, but doesn't explain why the diagnostics are synchronically ranked or why partitive \mention{Y}-licensing resists prosodic rescue more strongly than existential pivots. A simple feature account predicts all-or-nothing specificity effects rather than partial amelioration. An arbitrary dialect account predicts local frame differences without systematic preservation or loss of institutional readings.

The projectibility-first account does more than accommodate existing data; it generates empirical commitments about where diagnostics should split and where they should continue to travel together.

\section{Case studies: Weak definites, generic definites, and narrative \mentionhead{this}}\label{sec:cases}

The explanatory power of the account is best demonstrated by long-standing puzzles. This section uses weak definites, generic definites, and narrative \mention{this} as stress tests for projectibility. Once a construction is identified as deital, the analysis should let us predict which morphosyntactic residue survives when canonical definiteness drops out \parencite{FrancisMichaelis2003}. These constructions don't show a perfectly clean syntax/semantics split. They show something more useful: the structural and semantic profiles are statistically coupled, partly independent, and leaky in predictable places.

\subsection{Weak definites}

So-called weak definite constructions like \mention{take the bus}, \mention{listen to the radio}, and \mention{go to the hospital} have troubled semantic theories for decades \parencite{carlsonsussman2005,aguilarguevarazwarts2010}. The puzzle is how \mention{the} appears without unique or familiar referents. Some semantic analyses invoke incorporation, kind reference, or minimal situations to preserve or refine a uniqueness-based account. Those mechanisms are important, but on their own they don't explain the full pattern: extreme sensitivity to modification, strict lexical restrictions, and enriched meanings (\mention{go to the hospital} = `seek medical treatment').

The projectibility-first account reframes the puzzle. Weak definites use deital morphology and inherit some distributional residue from ordinary \mention{the}-marked NPs, but they don't uniformly satisfy every diagnostic. That mixed behaviour is central. They lack the definiteness-marking function that typically accompanies \mention{the}, but they behave unlike ordinary indefinites: they're pragmatically conventionalized frames in which deital form and token-level definiteness semantics have partly separated.

They don't contribute token-level uniqueness or familiarity. \mention{Take the bus} doesn't presuppose a unique, identifiable bus; it calls up a conventional activity type. The morphosyntactic frame (verb + \mention{the} + noun in this construction) is stable and deital in form, but the semantic contribution is non-unique reference. This form--meaning split has a permeable boundary between grammar and interpretation.

The frame does encode stereotypicality and conventionalization, which creates a weak form of identifiability: hearers recognize the activity type even when the token referent is unspecified. But this falls short of the definiteness profile's core properties (token uniqueness and anaphoric recoverability). Weak definites occupy the gap: morphosyntactically deital, semantically indefinite, pragmatically conventional.

The pattern can be described in terms of grammaticalization and constructional entrenchment, but that description has limits. Weak definites can be analysed as entrenched constructional frames in which \mention{the} persists because of institutional or habitual association while token-level definiteness is weakened or absent. Here the distributional residue (appears with \mention{the}, resists modification, allows only certain verb--noun pairs) is what remains of a construction after its ordinary token-level definiteness has weakened. This doesn't explain the full lexical inventory: why \mention{take the bus} and \mention{listen to the radio}, but not every institutionally associated noun. The account predicts selective residue once a weak-definite frame is identified; it doesn't derive every frame that English happens to contain.

Weak definites don't pattern uniformly across all diagnostics. In partitive constructions, \mention{two of the buses} is available only on an ordinary referential reading; \mention{two of the bus} is odd on the weak institutional reading. This follows from the diagnostic ranking: partitive \mention{Y}-licensing is the hardest syntactic test and doesn't simply yield to constructional pressure. The weak reading is available only when the frame itself licenses it. Partitives require a referentially established domain; weak definites provide institutional or activity-type identifiability. The mismatch blocks the weak reading.

\subsection{Generic definites}

So-called generic definites like \mention{The lion is noble} or \mention{The computer changed the world} pose a similar challenge. The definite singular is used for kind reference, not individual reference \parencite{carlson1977,krifka2004}. Uniqueness holds at the kind level, but the standard application of definiteness theory targets individuals.

Again, the projectibility-first account provides clarity, though generic definites are a compatibility case rather than the strongest independent argument. A semantic account can treat \mention{the lion} as unique or maximal at the kind level. The deitality claim is different: whatever the correct kind semantics, the NP still uses the deital morphosyntactic profile. Its behaviour belongs with the residue question: which distributional properties remain when individual-level definiteness isn't at issue?

English allows multiple strategies for genericity: bare plurals (\mention{Lions are noble}), indefinite singulars (\mention{A lion is noble}), and definite singulars (\mention{The lion is noble}). The choice between them is conditioned by subtle semantic and pragmatic factors. Definite singulars are preferred for well-established kinds, particularly species or inventions conceived as singular abstract individuals.

What the definite singular provides is the morphosyntactic profile of deitality. The semantics is kind reference, which operates at a different level from individual reference. Individual-level uniqueness and identifiability don't apply in the usual way; kind-level properties do. The syntax generates a deital phrase; semantics assigns it a generic interpretation. This supports the distinction, but it doesn't by itself force the distinction against every kind-based semantic account.

This point shouldn't be treated as independent deital residue. Generic definites resist existential \mention{there} because the construction is geared to episodic, stage-level predication, and kind-level subjects are semantically incompatible. Bare-plural generics show the same pressure. The deitality claim is narrower: a definite singular generic retains deital morphosyntax while its semantic interpretation operates at the kind level.

\subsection{Indefinite \mentionhead{this} as constructional recruitment}

Colloquial English productively uses \mention{this} with discourse-new referents in presentational and narrative contexts \parencite{prince1981,GernsbacherShroyer1989}:
\ea 
\ea \textit{There was this guy at the door...}
\ex \textit{So I'm walking home when this dog starts following me...}
\z
\z
\noindent This pattern is frequent in speech and systematic across speakers, so it isn't performance error or dialectal idiosyncrasy. But it strains the alignment between deitality and definiteness: demonstrative morphology appears without discourse-old or unique referents.

The account explains this as \term{constructional recruitment}. Certain narrative-presentation constructions recruit deital morphology for specificity, subjective immediacy, and cataphoric accessibility rather than ordinary definiteness. Demonstrative morphology is a motivated recruit here, not an arbitrary one: demonstratives already direct attention and encode speaker-anchored salience, which is just the material narrative presentation needs for vividness and expected continuation. The construction selects a demonstrative-derived form while introducing a discourse-new referent, exactly the kind of form--meaning mismatch the framework predicts when morphosyntactic frames achieve constructional autonomy (§\ref{subsec:autonomy}).

If narrative \mention{this} reflects constructional recruitment rather than a separate lexical item, we predict asymmetric diagnostic behaviour:

\begin{itemize}
    \item Existential \mention{there}. The narrative construction licenses the combination directly. \mention{There was this guy} is acceptable because the construction packages discourse-new reference with narrative prosody and cataphoric accessibility. The contrast with \mention{There was a guy at the door} isn't basic grammaticality; the latter is an ordinary narrative opening. The predicted contrast is discourse effect: \mention{this} should make the referent more vivid, noteworthy, and likely to be continued.

    \item Partitive \mention{Y}. The construction doesn't alter partitive \mention{Y}-licensing. A singular narrative \mention{this guy} can't supply a plural domain for \mention{two of}; plural demonstrative forms such as \mention{these guys} can serve as partitive complements only when they establish a referential group. Narrative framing doesn't by itself create that domain.

    \item Identificational hosting. Narrative \mention{this} should show intermediate acceptability, better than bare \mention{a} (because demonstrative morphology provides the salience hosting requires) but worse than canonical deital uses (because the referent lacks discourse history). This matches the specificity-forcing prediction in §\ref{sec:specific}.
\end{itemize}

Constructional autonomy (§\ref{subsec:autonomy}) explains how established morphosyntactic frames can be recruited for new functions, but it doesn't license every mismatch equally. English narrative presentation recruits demonstrative morphology for specificity, immediacy, and expected continuation while introducing a discourse-new referent. That recruitment is plausible because demonstratives already package attention direction.

The diagnostics needn't move together: the construction licenses the \mention{there} + demonstrative combination and improves discourse accessibility, but it doesn't automatically alter partitive domain requirements or the narrow \mention{one}-substitution residue. This selective pattern confirms that deitality is a profile of partly independent distributional requirements, which is why the diagnostics can move separately.

The sources of stability for narrative \mention{this} remain open. Three hypotheses would make the constructional account more specific:
\begin{itemize}
\item Acquisition packages the narrative frame with demonstrative morphology, so children don't have to decompose it into \enquote{indefinite this.}
\item Register and genre conventionalize the pattern in oral narrative, maintaining it across generations despite prescriptive pressure.
\item Processing supports the construction because its prosodic packaging (narrative contour) signals the non-canonical semantics.
\end{itemize}

The constructional account makes testable predictions that distinguish it from a polysemy analysis. Remove the narrative frame, and the discourse effect should weaken: \mention{This guy walked in} can be acceptable as a colloquial specific indefinite, but it should be less strongly cataphoric than \mention{There was this guy who walked in} in a story-opening frame. The construction should also differentiate determinatives by discourse effect rather than by simple grammaticality. \mention{There was a guy at the door...} should be acceptable, but \mention{There was this guy at the door...} should increase vividness, noteworthiness, and expected continuation of the referent. Finally, acquisition should show frame-dependence: children should initially restrict narrative \mention{this} to learned story-opening and presentational constructions before generalizing. They shouldn't show across-the-board \enquote{indefinite this} independent of constructional context.

If these predictions fail (if \mention{this} behaves exactly like \mention{a} in narrative frames, if it shows no special cataphoric effect, or if children use indefinite \mention{this} freely without constructional anchoring), the constructional account is falsified and a polysemy analysis becomes necessary.

\section{Edge cases and objections}\label{sec:objections}

\subsection{Proper names}

Proper names are semantically definite (they refer to unique, identifiable individuals), but they typically lack deital morphology in English \parencite{lyons1999}. The reason isn't that singular names are exempt from deital marking. It's that bare singular names sit outside the relevant part of the determiner system. In \CGEL{} terms, \mention{John} is a proper noun functioning as a full NP. It secures reference lexically, so the NP doesn't need a determiner slot whose filling could carry a deital signal.

The demonstrative-to-article path supplies the background for this prediction \parencite{greenberg1978,diessel1999}. Demonstrative-derived markers should spread most naturally where reference has to be secured compositionally: ordinary definite descriptions, anaphoric NPs, and related determiner-function constructions. Bare singular names are where overt identifiability marking does the least work. Modern Greek commonly uses an overt definite marker with personal names in ordinary argument positions, for example \olang{o Stavros} `the Stavros', though its use is conditioned by case, construction, and register \parencite{HoltonEtAl2012}. English bare \mention{John} and Greek \olang{o Stavros} are positions in a diachronic and constructional space, not counterexamples to the deitality--definiteness correlation.

Within English, deital marking reasserts itself exactly where names enter ordinary NP determination. Pluralization forces property denotation: \mention{Smiths} denotes something like \enquote{members of the Smith family}, not a single rigid individual. Once the name has been coerced into a plural nominal, the productive bare-plural pattern would suggest an indefinite or generic reading. The deital determiner marks the intended identifiable plural domain: \mention{the Smiths}, \mention{the Johnsons}. This isn't an ad hoc override. It's ordinary deital determination after coercion into a construction where determination has work to do.

The same mechanism covers nearby cases. \mention{A Susan arrived} coerces the name into a common-noun-like description, roughly `some person named Susan'. \mention{The young Einstein} uses a name with descriptive content and deital marking. Lexicalized \mention{the}-names such as \mention{The Boring Company} package the article into the name's conventional form. A bad string such as \mention{*two of John} doesn't show non-deitality by itself; it's bad because a singular name denotes an individual rather than a plural domain. More revealing are cases where the name is coerced, pluralized, coordinated, or lexicalized, since those cases show when names do and don't enter the deital profile.

This gives a projection target. Deital behaviour should strengthen as names move from bare lexical reference into ordinary NP determination: bare singular names should show the least deital behaviour; modified names, plural family names, lexicalized \mention{the}-names, and name-derived common nouns should show more of it when the intended reading is identifiable. The account doesn't predict where a given language sits on the cline. That's contingent history. It predicts the shape of the variation: names can be definite without deital marking, require deital marking when constructional cues would otherwise mislead, or commonly appear with deital marking in ordinary argument positions.

\subsection{Interrogatives and relatives}

Interrogative and relative determinatives require special treatment. Most deitality diagnostics presuppose declarative structure. Interrogatives have obligatorily fronted phrases; relatives have distinct clause-internal syntax.

In partitive constructions under standard interrogative force, interrogatives can't serve as \mention{Y}: \mention{*How many of which students passed?} is degraded compared to \mention{How many of the students passed?} Echo questions allow it (\mention{Two of WHICH students?}), but the construction is marked, with distinct prosody and pragmatics. When interrogatives appear in \mention{X} position (\mention{which of the students}, \mention{how many of these books}), they enforce definite-like \mention{Y} just as other quantifiers in \mention{X} do, but the contrast reflects the partitive construction's selectional requirements rather than the interrogative's own deital status. Interrogatives also license bare \mention{one} (\mention{Which one?}), unlike deital determinatives. Under neutral conditions, then, the visible diagnostics don't establish a deital profile; I mark the rows Provisional in the inventory.

For relatives, I follow \emph{The Cambridge Grammar of the English Language} in treating \mention{that} as a subordinator, not a determinative \parencite{HuddlestonPullum2002}. It marks subordination and doesn't realize an argument position. \mention{Which} is a determinative when it appears with an overt head (\mention{in which chapter}). Relative \mention{which} also shows definiteness without ordinary deital marking: \mention{It may rain, in which case we'll leave early} is anaphoric and definite in interpretation, while the independent-clause paraphrase recruits a deital demonstrative, \mention{It may rain. In that case, we'll leave early}. In fused-head constructions (\mention{I read which(ever) you recommended}), \mention{which} patterns like other fused forms, sometimes deital, sometimes not, depending on the broader structure.

\subsection{Exclamatives}

Exclamative phrases like \mention{What a mess!} or \mention{So many people!} have clause-type-specific syntax. Most diagnostics don't apply cleanly. Where exclamative phrases surface in declarative contexts, they pattern non-deitally: \mention{There was such a mess} is fine; \mention{??Such a book, which I bought, is rare} is degraded.

Exclamative force (expressing extreme degree) is orthogonal to the deitality profile. When exclamative determinatives appear outside their clause type, they behave like other degree-modified quantifiers: acceptable after \mention{there}, poor hosts for identificational constructions.

\subsection{Circularity and collapse objections}

The most serious objection is circularity: the diagnostics define the category they purport to identify. If deitality is characterized by resistance to existential \mention{there}, partitive \mention{Y}-licensing, etc., and those same properties are used to identify deital determinatives, haven't we just described the same facts twice?

The reply begins by separating three moments in the argument. Diagnostics identify a synchronic profile. Diachrony and acquisition explain why that profile isn't accidental. Projection tests ask whether locating an item or construction in the profile licenses further inferences. Identification doesn't have to be independent of every diagnostic property; anti-essentialist categories are often identified by clustered symptoms. The requirement is that the profile support inferences beyond mere relabelling.

Two safeguards keep the argument from merely naming its own evidence. First, the diagnostics are independently motivated: existential \mention{there}, partitives, and identificational constructions are well-studied phenomena with their own syntactic and semantic analyses \parencite{milsark1977,Barker1998,Lambrecht1994}. I draw on these established tests to reveal a profile. Second, the strongest sources and stabilizers are independently attested. The cross-linguistic record documents grammaticalization pathways from demonstratives to definite markers \parencite{diessel1999}. Acquisitional overgeneralization of \mention{the} is an empirical finding \parencite{maratsos1976}. The account doesn't infer every stabilizer from the data it's meant to explain. It derives predictions from independently motivated diachronic and acquisitional processes.

The deeper worry is collapse: perhaps the diagnostics aren't a profile at all but one semantic property (presuppositionality, referentiality, specificity, or familiarity) measured four ways. The answer is dissociation under manipulation. Prosody can rescue existential pivots but not partitive \mention{Y} complements. Specificity can improve identificational hosting without changing partitive \mention{Y}-licensing. Weak definites preserve deital \mention{the}-morphology but lose the weak reading under partitive \mention{Y}-licensing pressure. Quantifiers such as \mention{each} and \mention{every} resist existential pivots while failing other diagnostics. These splits would be surprising if the diagnostics were just one hidden semantic feature. They're expected if deitality is a ranked profile whose components are related but not identical.

A more sophisticated collapse objection replaces one semantic feature with several: existential \mention{there} tracks information-structural newness, partitives track domain establishment, hosting tracks specificity, and the observed clustering is just the ordinary correlation among those semantic and pragmatic properties. This version is harder to dismiss, and one part is right. The diagnostics are semantically leaky. A semantic-bundle account can also mimic some of the manipulation results: it can say that domain establishment is prosody-insensitive while presentational newness is prosody-sensitive. The prosody split alone doesn't settle the issue.

The projectibility-first account earns its keep only where it predicts effects that a semantic bundle doesn't get without extra assumptions. Two tests matter most. The first is acquisitional. If partitive \mention{Y}-licensing, existential pivot resistance, and hosting are learned as unrelated semantic filters, there's no clear reason to expect the ordering predicted above: earlier control of the hard partitive restriction and later adult-like control of neutral-prosody existential pivots. If children instead learn all three according to independent semantic complexity, the projectibility-first account loses support.

The second test concerns weak definites in existential contexts. On the projectibility-first account, an institutional weak definite shouldn't freely behave like an ordinary indefinite pivot just because its token-level semantics is non-unique. A sentence such as \mention{There's the bus} is acceptable on deictic, arrival, or list readings; the discriminating question is whether the weak institutional reading survives a neutral expletive-\mention{there} frame as readily as ordinary indefinite pivots do. If it does, the semantic-bundle account gains ground. If the institutional reading resists neutral pivothood while deictic or presentational readings are available, that's deital residue not explained by indefinite semantics alone.

The same historical forms also recur across the distributional profile: demonstrative-derived determinatives, including English \mention{the}, \mention{this}, \mention{that}, \mention{these}, and \mention{those}, carry a shared formal lineage described above. The /ð/ pattern isn't a diagnostic, but it's a trace of form history that a semantic-only account treats as accidental. The positive case is cumulative: lineage, acquisition, weak-definite residue, and cross-level manipulation have to pattern together. Any one of them can be explained away locally; the profile claim says they should travel together often enough to support further inference.

\section{Conclusion}\label{sec:conclusion}

The definiteness problem becomes clearer once we stop asking one profile to do the work of two. Deitality and definiteness are strongly correlated, but they support different inferences. Deitality projects from form to distribution: from \mention{the}, demonstratives, and related determiner-function constructions to expectations about partitives, existential pivots, identificational hosting, and the narrow \mention{one}-substitution residue. Definiteness projects from interpretation to discourse management: from identifiability, uniqueness, and recoverability to expectations about anaphora, common ground, and topic continuity.

The mismatch cases show why the separation matters. Weak definites, generic definites, narrative \mention{this}, and proper names show that morphological marking and referent identifiability don't align perfectly \parencite{FrancisMichaelis2003}. Two illustrate the point most sharply: weak definites show deital morphology without ordinary token-level definiteness semantics, while proper names show definiteness semantics without ordinary deital marking. Generic definites and narrative \mention{this} extend the same logic to kind reference and constructional recruitment. Together, the mismatch cases let us see which parts of the two profiles remain coupled and which parts can move independently.

The natural-kind apparatus adds one constraint: a category earns its place by supporting further inference, not by supplying an essence or a strict corrective-control system. The analysis is limited in the same way. A semantic-bundle account can mimic some of the dissociations, so the strongest tests are the ones that separate the accounts, especially acquisition order and weak-definite behaviour in neutral existential frames.

The predictions run from utterance-level manipulation to community-level variation and developmental change. Prosody should affect existential pivots but not partitive \mention{Y}-licensing. Specificity should improve indefinites in identificational environments without producing deital baselines. Licensed institutional frames should preserve institutional readings where demonstratives and genitive NPs tend to force referential ones. Broader dialect data should show a layered pattern: article realization, article presence, reinforced demonstratives, and demonstrative paradigms may vary across dialect grammars, while a variety that allows weak indefinites as ordinary partitive \mention{Y} complements would mark a deeper grammar difference. Children should show earlier control of partitive \mention{Y}-licensing than adult-like neutral-prosody \mention{there}-resistance. These are empirical commitments, not just redescriptions.

Two broader lessons follow. Grammatical categories can be profile-like rather than essence-like. Sometimes the right move isn't a better definition of the inherited label but a split into profiles with different diagnostics, stabilizers, and projection targets. Aspect, mood, transitivity, verb, and adjective are possible candidates for the same treatment. The maintenance story beyond grammaticalization and acquisition is deliberately scoped as a research programme: the result here is the projectible profile and its predictions, not a complete causal theory of every stabilizer.

Deitality and definiteness remain strongly correlated because grammaticalization, acquisition, and discourse pressure keep drawing them together, but the categories are distinct. Kept distinct, they explain why diagnostics travel together, why mismatches preserve deital residue, and why prosody, construction, dialect, and acquisition perturb different parts of the system. The point is to let each category do its own work.

\clearpage
\appendix
\section{Provisional Inventory of English Determinatives and Related Forms}\label{sec:inventory}

The table below is a compact orientation device; it isn't a defended taxonomy of every English determinative. The cells are author judgments intended to make the profile auditable; only the core contrasts developed in the main text are used as evidence in the argument. Rows marked Provisional, especially fused relatives, exclamatives, compound forms, and personal determinatives, mark empirical terrain for later judgment and corpus work.

\begingroup
\footnotesize
\begin{longtable}{@{}lcccl@{}}
\caption{Provisional inventory of English determinatives and related forms.}\label{tab:inventory}\\
\toprule
\textbf{Item or construction} & \textbf{\mention{there}-resist.} & \textbf{Domain-\mention{Y}} & \textbf{Hosting} & \textbf{Status} \\
\midrule
\endfirsthead
\caption[]{Provisional inventory of English determinatives and related forms (continued).}\\
\toprule
\textbf{Item or construction} & \textbf{\mention{there}-resist.} & \textbf{Domain-\mention{Y}} & \textbf{Hosting} & \textbf{Status} \\
\midrule
\endhead
\midrule
\multicolumn{5}{r}{\emph{Continued on next page}}\\
\endfoot
\bottomrule
\endlastfoot
\multicolumn{5}{l}{\emph{Core deital items and constructions}} \\
\mention{the} & \checkmark & \checkmark & \checkmark & \Deital{} \\
\mention{this} & \checkmark & \checkmark & \checkmark & \Deital{} \\
\mention{that} & \checkmark & \checkmark & \checkmark & \Deital{} \\
\mention{these} & \checkmark & \checkmark & \checkmark & \Deital{} \\
\mention{those} & \checkmark & \checkmark & \checkmark & \Deital{} \\
personal genitives (\mention{my}, \mention{your}, etc.) & \checkmark & \checkmark & \checkmark & \Deital{} \\
genitive NPs in determiner function (\mention{John's}) & \checkmark & \checkmark & \checkmark & \Deital{} \\
\addlinespace
\multicolumn{5}{l}{\emph{Universal/distributive quantifiers} (mixed)} \\
\mention{each} & \checkmark & -- & -- & Mixed \\
\mention{every} & \checkmark & -- & -- & Mixed \\
\mention{everybody/everyone}$^a$ & \checkmark & -- & ? & Provisional \\
\mention{everything}$^a$ & \checkmark & -- & ? & Provisional \\
\mention{everywhere}$^a$ & n/a & n/a & ? & Provisional \\
\mention{all} & --$^b$ & -- & -- & Mixed \\
\mention{both} & --$^b$ & -- & -- & Mixed \\
\addlinespace
\multicolumn{5}{l}{\emph{Alternative/negative quantifiers}} \\
\mention{either} & \checkmark & -- & -- & Mixed \\
\mention{neither} & \checkmark & -- & -- & Mixed \\
\mention{no} & -- & -- & -- & \Nondeital{} \\
\mention{nobody/no one}$^a$ & -- & -- & -- & \Nondeital{} \\
\mention{nothing}$^a$ & -- & -- & -- & \Nondeital{} \\
\mention{nowhere}$^a$ & n/a & n/a & -- & Provisional \\
\mention{none} & -- & -- & -- & \Nondeital{} \\
\addlinespace
\multicolumn{5}{l}{\emph{Indefinites and weak quantifiers}} \\
\mention{a/an} & -- & -- & -- & \Nondeital{} \\
\mention{a certain} & -- & -- & -- & \Nondeital{} \\
\mention{a few} & -- & -- & -- & \Nondeital{} \\
\mention{a little} & -- & -- & -- & \Nondeital{} \\
\mention{another} & -- & -- & -- & \Nondeital{} \\
\mention{some} & -- & -- & -- & \Nondeital{} \\
\mention{somebody/someone}$^a$ & -- & -- & -- & \Nondeital{} \\
\mention{something}$^a$ & -- & -- & -- & \Nondeital{} \\
\mention{somewhere}$^a$ & n/a & n/a & -- & Provisional \\
\mention{any} & -- & -- & -- & \Nondeital{} \\
\mention{anybody/anyone}$^a$ & -- & -- & -- & \Nondeital{} \\
\mention{anything}$^a$ & -- & -- & -- & \Nondeital{} \\
\mention{anywhere}$^a$ & n/a & n/a & -- & Provisional \\
\addlinespace
\multicolumn{5}{l}{\emph{Degree/quantity expressions}} \\
\mention{many} & -- & -- & -- & \Nondeital{} \\
\mention{much} & -- & -- & -- & \Nondeital{} \\
\mention{few} & -- & -- & -- & \Nondeital{} \\
\mention{fewer} & -- & -- & -- & \Nondeital{} \\
\mention{little} & -- & -- & -- & \Nondeital{} \\
\mention{less} & -- & -- & -- & \Nondeital{} \\
\mention{more} & -- & -- & -- & \Nondeital{} \\
\mention{most} (proportional) & -- & -- & -- & \Nondeital{} \\
\mention{several} & -- & -- & -- & \Nondeital{} \\
\mention{certain} (plural) & -- & -- & -- & \Nondeital{} \\
\mention{various} & -- & -- & -- & \Nondeital{} \\
\mention{multiple} & -- & -- & -- & \Nondeital{} \\
\mention{enough} & -- & -- & -- & \Nondeital{} \\
\mention{sufficient} & -- & -- & -- & \Nondeital{} \\
\mention{umpteen} & -- & -- & -- & \Nondeital{} \\
\addlinespace
\multicolumn{5}{l}{\emph{Superlatives and ordinals} (deital in referential frames)} \\
(\mention{the}) \mention{most\upslash fewest\upslash least}$^c$ & \checkmark & -- & \checkmark & \Deital{} \\
\mention{last}$^c$ & \checkmark & n/a & \checkmark & \Deital{} \\
\mention{next}$^c$ & \checkmark & n/a & \checkmark & \Deital{} \\
\mention{said} & \checkmark & \checkmark & \checkmark & \Deital{} \\
\addlinespace
\multicolumn{5}{l}{\emph{Numerals}$^d$} \\
\mention{zero} & -- & -- & -- & \Nondeital{} \\
\mention{one} (numeral) & -- & -- & -- & \Nondeital{} \\
\mention{two, three, four...} & -- & -- & -- & \Nondeital{} \\
\mention{ten, twenty...ninety} & -- & -- & -- & \Nondeital{} \\
\mention{hundred, thousand, million} & -- & -- & -- & \Nondeital{} \\
\mention{once, twice, thrice}$^d$ & -- & -- & -- & \Nondeital{} \\
\addlinespace
\multicolumn{5}{l}{\emph{Interrogatives} (clause-type dependent)$^e$} \\
\mention{what} & n/a & -- & n/a & Provisional \\
\mention{which} & n/a & -- & n/a & Provisional \\
\mention{whose} & n/a & -- & n/a & Provisional \\
\mention{how many/much} & n/a & -- & n/a & Provisional \\
\mention{what/which kind of} & n/a & -- & n/a & Provisional \\
\addlinespace
\multicolumn{5}{l}{\emph{Relatives and fused relatives}$^f$} \\
\mention{which} (with overt head) & n/a & -- & n/a & Provisional \\
\mention{what} (fused relative) & \checkmark & \checkmark & \checkmark & Provisional \\
\mention{whatever} (ignorance) & \checkmark & \checkmark & \checkmark & Provisional \\
\mention{whatever} (free-choice) & \checkmark & -- & -- & Mixed \\
\mention{whichever} (ignorance) & \checkmark & \checkmark & \checkmark & Provisional \\
\mention{whichever} (free-choice) & \checkmark & -- & -- & Mixed \\
\addlinespace
\multicolumn{5}{l}{\emph{Exclamatives} (clause-type dependent)$^g$} \\
\mention{what} (exclamative) & -- & -- & -- & \Nondeital{} \\
\mention{what a} & -- & -- & -- & \Nondeital{} \\
\mention{such} & -- & -- & -- & \Nondeital{} \\
\mention{such a} & -- & -- & -- & \Nondeital{} \\
\mention{so many/much} & -- & -- & -- & \Nondeital{} \\
\mention{how many/much} (excl.) & -- & -- & -- & \Nondeital{} \\
\addlinespace
\multicolumn{5}{l}{\emph{Personal determinatives}$^h$} \\
\mention{us} & n/a & ? & n/a & Provisional \\
\mention{we} & n/a & ? & n/a & Provisional \\
\mention{you} & n/a & ? & n/a & Provisional \\
\end{longtable}
\endgroup

\begingroup
\footnotesize
\raggedright
\noindent\checkmark{} = shows the \textsc{deital} pattern named by the column; -- = fails the pattern; n/a = diagnostic undefined for the construction type; ? = untested or confounded in the relevant construction. \mention{There}-resist. means \emph{resists} post-\mention{there} pivothood under neutral prosody (i.e., exhibits the definiteness effect), not that the item occurs after \mention{there} \parencite{milsark1977,McNally2011}. Domain-\mention{Y} means that the item can supply the domain-establishing \mention{Y} complement in a true partitive. The narrow \mention{one}-substitution residue is discussed in §\ref{sec:diagnostics} but isn't used as an inventory column. For status, \textsc{deital} = shows convergence across the ranked diagnostics (prototypical profile members); \textsc{non-deital} = fails the ranked diagnostics; Mixed = shows partial profile convergence rather than a settled boundary verdict; Provisional = included as a target for later validation rather than used as evidence in the paper.

\noindent$^a$The compound rows follow a \CGEL-style category/function lens; readers who treat these as pronouns or adverbs should read them as fused or extended determinative-related expressions rather than core determinative lexemes \parencite{HuddlestonPullum2002}. $^b$Coordination rescues \mention{there}-pivot. $^c$These rows concern referential determiner-function uses, not proportional quantifier use; bare \mention{last} and \mention{next} don't independently supply a full Domain-\mention{Y} complement, so the diagnostic is undefined rather than failed. $^d$See \PrepositionCitation{} for the inclusion of numeral-derived forms such as \mention{once}, \mention{twice}, and \mention{thrice}. $^e$Interrogatives are fronted in questions; diagnostics are largely undefined. $^f$Fused relatives pattern varies by reading (ignorance-definite vs. free-choice). $^g$Exclamatives in declarative contexts pattern non-deitally; frame-initial uses are undefined. $^h$Personal determinatives are limited to constructions such as \mention{us linguists} and require separate treatment; the partitive contrast \mention{two of us/you} vs. \mention{*two of we} is confounded by case selection by \mention{of}.
\par
\endgroup

\section*{Declarations}

\paragraph{Funding.} No external funding was received for this work.

\paragraph{Competing interests.} The author declares no competing interests.

\paragraph{AI assistance.} During manuscript preparation and revision, primarily in June 2026, the author used Claude Sonnet 4.5 extensively for drafting and editing support, with some input from ChatGPT 5 and Gemini 2.5 Pro. The author also used Claude Opus 4.8, Claude Fable 5, and GPT-5.5 Pro for review feedback, and OpenAI Codex (GPT-5), accessed through the Codex CLI, for revision support. These tools were used through their standard Anthropic, OpenAI, and Google interfaces. The author reviewed and edited all text, verified sources and references, and takes full responsibility for all claims and conclusions.

\clearpage
\printbibliography

\end{document}
